% Compared with arXiv:2512.15383v1; AUTHOR-CHECK notes remain author decisions.
\documentclass[10pt]{article}
\usepackage[margin=0.72in]{geometry}
\usepackage[T1]{fontenc}
\usepackage{lmodern,amsmath,enumitem,xcolor,microtype}
\usepackage{xr-hyper}
\usepackage[colorlinks=true,linkcolor=black,urlcolor=blue,bookmarks=false]{hyperref}
\externaldocument[M-]{main}[main.pdf]
\setlength{\parindent}{0pt}
\setlength{\parskip}{0.35em}
\setlength{\emergencystretch}{2em}
\setlist[itemize]{leftmargin=1.25em,itemsep=0.55em,topsep=0.3em}
\newcommand{\mssec}[1]{Section~\ref*{M-#1}}
\newcommand{\msapp}[1]{Appendix~\ref*{M-#1}}
\newcommand{\checkauthor}[2]{\textcolor{red!65!black}{\textbf{AUTHOR-CHECK #1:} #2}}
\begin{document}
\begin{center}
{\Large\bfseries Summary of Revisions}\par
\vspace{0.3em}
{\bfseries Interpretable Multivariate Conformal Prediction\par
with Fast Transductive Standardization}\par
JMLR-25-3153-1
\end{center}

This summary compares the current draft with the original arXiv v1
submission of December 17, 2025. It describes changes present in the draft.

\begin{itemize}
\item \textbf{Abstract.} No changes.

\item \textbf{Section 1: introduction and related work.} The motivation for
interpretable simultaneous intervals and the Bonferroni comparison are
retained. Figure 1 now shows illustrative marginal coverages of 92\% and
93\% alongside 90\% joint coverage, making that distinction visible
immediately. Added a reference Tumu et al. for rectangular shape templates and clarified that Sampson's method has two versions and does not always require additional sample splits.

\item \textbf{Section 2: setup and motivating constructions.} The revised
setup states joint exchangeability of calibration and test observations
and explains that coverage is marginal over the data but simultaneous
across outputs. Residual definitions now discuss capping and ties. We now use a generic
location/scale parameterization of the prediction set to separate the population oracle from its
practical approximations with only different plug-in estimators. This makes the presentation lighter.

\item \textbf{Section 3: methodology.} Similarly, we now write the most general form of an estimated prediction set in the beginning, so that everything we develop in this section is a plug-in approximation of the same oracle. Figure 2 labels
the residual boundaries; multi-indices are consistently bold. The
real-data diagnostic reference now points to the appendix. These changes
make the construction easier to follow.

\item \textbf{Section 4: numerical experiments.}
\textbf{Twelve experimental studies along with 22 new plots have been added in total: ten simulation
studies and two real-data diagnostic studies.} This count spans Section 4
and Appendix C. 
\begin{enumerate}
    \item \textbf{E01: Dependent heterogeneous Gaussian noise (\mssec{sec:absolute-residual-experiments}).}
Tests whether coordinate adaptation remains useful when outputs are correlated. Joint coverage, volume, runtime, and coordinate lengths distinguish scale adaptation from the smaller but under-covering copula regions.

\item \textbf{E02: Partial heteroskedasticity (\mssec{sec:absolute-residual-experiments}).}
Makes only five of ten coordinates input-dependent while preserving marginal second moments. 

\item \textbf{E03: Quantile regression experiments (\mssec{sec:cqr-experiments}).}
Tests quantile-regression residuals using capped transformation and adds the requested CQHR baseline. Outcome-space volumes and lengths show design-dependent gains of capped TSCP while identifying its inability to shrink the fitted base interval.
\end{enumerate}

\item \textbf{Section 4.4: real-data presentation.} The original results
are retained in an expanded table. A new plot of existing timing records
makes wall-clock cost visible in the body. The Point CHR infinite volume on \texttt{stock} is explained
by the small final calibration split. Coordinate and search diagnostics
are presented in Appendix C.

\item \textbf{Section 5: discussion.} No changes.
\end{itemize}

\begin{itemize}
\item \textbf{Appendix A: algorithmic details and complexity.} Fixed some typos.

\item \textbf{Appendix B: proofs.} No changes besides updating the notation according to the body.

\item \textbf{Appendix C: supplementary experiments.} 

\begin{enumerate}
    \item \textbf{E04: Dependence-strength test (\msapp{app:dependence-sensitivity}).}
Varying correlation at fixed marginal scales and shows stable TSCP coverage and decreasing volume as dependence strengthens.

\item \textbf{E05: Coordinate-heterogeneity test(\msapp{app:heterogeneity-sensitivity}).}
Identifies when coordinate rescaling is useful by varying the largest-to-smallest noise-scale ratio. The common unscaled threshold becomes less efficient as heterogeneity increases, while remaining competitive at equal scales.

\item \textbf{E06: Target-coverage test(\msapp{app:alpha-small-sample}).}
Checks that the comparison is not specific to a 90 percent target. Targets of 80, 90, and 95 percent show the coverage-volume trade-off and the relative behavior of the methods.

\item \textbf{E07: Exchangeable contamination (\msapp{app:contamination}).}
Tests the sensitivity of mean and standard deviation scaling to outliers. Increasing row-wise contamination shows that near-nominal coverage can coexist with severe volume inflation.

\item \textbf{E08: CQR base-interval sensitivity (\msapp{app:cqr-sensitivity}).}
Checks dependence on the width of the fitted starting interval. The base-miscoverage testreveals CQHR's sensitivity to fitted widths and puts the main CQR comparison in context.

\item \textbf{E09: Shifted CQR scores (\msapp{app:cqr-sensitivity}).}
Examines additive shifts as an alternative to capping signed scores. Constants 100, 300, and 1000 give nearly identical recorded rectangles, equal to shifted TSCP-GWC within each run.

\item \textbf{E10: Rectangular shape-template baseline (\msapp{app:shape-template}).}
Adds the requested comparison with the Tumu et al. family. The two-dimensional KDE/mean-shift study compares coverage, normalized volume, and runtime for a specific rectangular template implementation, not every multimodal design.

\item \textbf{E11: Real-data coordinate diagnostics (\msapp{app:real-coordinate-audit}).}
Use now coordinate-wise length and marginal coverages comparisons to show where volume gains do and do not shorten individual intervals.

\item \textbf{E12: Real-data search and fallback diagnostics (\msapp{app:real-data-diagnostics}).}
Measures how often backward search or fallback occurs at different miscoverage levels. No backward search or fallback occurs at miscoverage 0.1 in the tested splits; larger miscoverage levels exercise both branches without demonstrating a worst-case scan.
\end{enumerate}

\end{itemize}
\end{document}
